\section{A certified obstruction to the tempting constant-four ramp bound}
\label{sec:revisit-ramp-four}
The first numerical experiments suggested
$\|\mD^{-1/2}(\vx_k-\vx_0^*)\|_\infty\leq4r_k$ during the ramp.
That particular constant is false. We first give certified finite witnesses;
\cref{sec:revisit-ramp-asymptotic} then rules out every graph-independent
constant by a separate asymptotic argument.

The original guess has an exact linear explanation. When $\alpha=\theta^2$
and the orthant constraint is omitted, the exponential schedule admits the
particular trajectory
\[
 \vx_k=\vx_0^*-4\eta r_k\vomega,\qquad
 \vz_k=\vx_0^*-2r_k\vomega.
\]
Substitution uses $\mQ\vomega=\alpha\vomega$ and
$\eta-\chi=\theta/2$. Its error factor $4\eta<4$ approaches four as
$\alpha$ decreases. This particular solution does not have the algorithm's
zero initial state and can have negative coordinates. It therefore controls
neither the homogeneous transient nor changes of the active constraints.

\begin{lemma}[Pairwise contraction of the exact step]
\label{lem:revisit-pairwise}
For two incoming states, apply the exact orthant recurrence with the same
source and regularizer. On state differences define
\[
 \mathcal D(\vd_x,\vd_z)
   =\left(\tfrac12\|\vd_x\|_{\mQ}^2
                   +\tfrac\mu2\|\vd_z\|_2^2\right)^{1/2}.
\]
Then $\mathcal D(\vd_x^+,\vd_z^+)\leq\sqrt\chi\,
\mathcal D(\vd_x,\vd_z)$.
\end{lemma}
\begin{proof}
Subtract the two raw recurrences. The source and regularizer cancel, so
their difference has gradient $\mQ\vd_y$.
Monotonicity of the orthant normal gives
\[
 (\vp_1-\vp_2)^\trans
       \{(\vq_1-\vp_1)-(\vq_2-\vp_2)\}\geq0.
\]
Apply the imported comparison identity to the quadratic
$\varphi(\vd)=\tfrac12\|\vd\|_{\mQ}^2$, comparator zero, and kinetic
output $\vp_1-\vp_2$. Taking square roots proves the assertion.
\end{proof}

\paragraph{An offline trajectory certificate.}
For this witness $\alpha=\theta^2=1/T^2$, $T=256$.
In degree-density coordinates, round both the kinetic and the next primal
state downward to a grid $h=2^{-80}$ after each dense exact step.
Also round $r$ downward to that grid, starting from the rounded $1/d_v$.
This is solely a verification procedure, separate from the local rounded
algorithm of \cref{sec:revisit-rounding}.

Its regularizer error is at most
$h/(1-\eta)=2h/\theta$. At a fixed incoming rounded state this perturbs
the raw kinetic density by at most $2h$, since $\alpha=\theta^2$.
Orthant projection and its floor consequently perturb the new kinetic
density by at most $3h$, and the primal density by at most $(1+3\theta)h$.
If $V$ is the full original graph volume, the perturbation in
$\mathcal D$ is at most $(1+3\theta)h\sqrt V$.
\Cref{lem:revisit-pairwise} and
$1-\sqrt\chi\geq\theta/2$ give, at every finite horizon,
\[
 \mathcal D(\vx_{\rm approx}-\vx_{\rm exact},
            \vz_{\rm approx}-\vz_{\rm exact})
 \leq5h\sqrt V/\theta.
\]
Thus the error in any primal degree density is at most
$5\sqrt2h\sqrt V/\theta^2$, which is less than the rational bound
$24\lceil\sqrt V\rceil T^2h$ used in the audit. Ambient volume is allowed
in this external verification certificate; it is not supplied to the
proposed local solver.

For the PPR reference, propose a dyadic vector $\vu$ from a numerical solve
and compute its residual exactly. Inverse positivity and
$\mQ\vomega=\alpha\vomega$ give the independent certificate
\[
 \|\mD^{-1/2}(\vu-\vx_0^*)\|_\infty
 \leq\|\mD^{-1/2}(\mQ\vu-\vb)\|_\infty/\alpha.
\]
The accuracy of the numerical solve is therefore not an assumption.

\begin{proposition}[Constant four fails]
\label{prop:revisit-four-fails}
Take a unit-weight clique on 128 vertices and attach a path of 384 new
vertices to one clique vertex. Seed a different clique vertex. For
$\alpha=1/65536$, $\theta=1/256$, $\eta=511/512$, and
$\rho=1/(1024\vol(\mathcal V))$, the exact ramp at step 5468 satisfies
\[
 4.01807<
 \frac{\|\mD^{-1/2}(\vx_{5468}-\vx_0^*)\|_\infty}{r_{5468}}
 <4.01809,
 \qquad r_{5468}=\frac1{127}\left(\frac{511}{512}\right)^{5468}>\rho.
\]
\end{proposition}
\begin{proof}[Certified finite verification]
The executable certificate is \texttt{ramp\_four\_counterexample.py} in the
note's proof-audit directory. Symmetry compresses the graph into the seed,
the other 126 non-connector clique vertices, the connector, and the 384
individual path vertices. Original degrees, class cardinalities, and
neighbor multiplicities are explicit; detailed balance and every row sum
are checked. The source and trajectory are constant on these classes.

The script evolves integer grid numerators for exactly 5468 steps. All
floors are integer divisions, so no trajectory arithmetic is floating point.
It then computes the exact rational PPR residual certificate and adds the
trajectory error bound above. Exact integer comparisons verify the two
displayed decimal rational bounds and that the ramp has not reached $\rho$.
This proves a finite counterexample, not an asymptotic lower bound.
\end{proof}

\subsection{Radial trees also refute monotonicity and PPR lower order}
Let $\mathcal T(b,s,L)$ be the rooted unit tree of depth $L$ in which a
vertex at level $j<L$ has $b$ children when $s$ divides $j$, and one child
otherwise. Seed the root. For all examples below use
$\rho=1/(1024\vol(\mathcal V))$; the displayed times occur strictly before
the ramp reaches $\rho$.

\begin{proposition}[Further certified ramp limitations]
\label{prop:revisit-radial-failures}
The following statements hold for the exact recurrence.
\begin{enumerate}
\item On $\mathcal T(10,3,12)$, which has 33,331 vertices, take
$\alpha=1/16384$ and $\theta=1/128$. The root degree density decreases
between steps 2249 and 2250 by an amount strictly between
$1.59\cdot10^{-11}$ and $1.61\cdot10^{-11}$.
\item On $\mathcal T(10,4,128)$, take $\alpha=1/4096$ and
$\theta=1/64$. At step 3037 the root density exceeds exact PPR by an
amount strictly between $4\cdot10^{-13}$ and $5\cdot10^{-13}$.
The minimum normalized PPR residual is strictly between
$-1.93\cdot10^{-14}$ and $-1.92\cdot10^{-14}$.
This tree has $1+4\sum_{j=1}^{32}10^j$ vertices.
\item On $\mathcal T(2,4,128)$, which has $8\cdot2^{32}-7$ vertices,
take $\alpha=1/4096$ and $\theta=1/64$. At step 3498,
\[
 6.82360<
 \frac{\|\mD^{-1/2}(\vx_{3498}-\vx_0^*)\|_\infty}{r_{3498}}
 <6.82361.
\]
\item On $\mathcal T(10,2,128)$, with
$1+2\sum_{j=1}^{64}10^j$ vertices, take $\alpha=1/4096$ and
$\theta=1/64$. At step 19129 the same error-to-$r$ ratio is strictly
between $2{,}700{,}097{,}454{,}229{,}948$ and
$2{,}700{,}097{,}454{,}229{,}950$.
\end{enumerate}
\end{proposition}
\begin{proof}[Certified finite verification]
The certificate \texttt{radial\_ramp\_counterexamples.py} uses one
equitable class per level. Class cardinalities are the successive products
of the branching factors; a nonroot vertex has one parent and its declared
number of children. The script checks row degrees and detailed balance,
so these are canonical finite unit graphs, not weighted quotient instances
substituted for the problem.

The integer trajectory uses grid $2^{-200}$ for the first three examples
and $2^{-700}$ for the fourth, with the same uniform error radius
$E=24\lceil\sqrt V\rceil T^2 h$ as above.
Exact rational tridiagonal elimination supplies the PPR reference; all of
its equations and its unit weighted mass are checked exactly. For a
one-step decrease use error radius $2E$; for an overshoot, the global
tracking norm, or a normalized residual use $E$ (the residual radius uses
$|\mQ|\vomega\leq\vomega$). Exact rational comparisons verify each
displayed bound and each condition $r_k>\rho$.
\end{proof}

The second example also disproves lower order to the current RPPR optimum,
because that optimum lies below exact PPR. These counterexamples justify
keeping the final accuracy-and-clipping phase and allowing signed residuals
and support excursions in the work proof. The finite witnesses alone do
not disprove all larger constants; that conclusion uses the separate proof
in \cref{sec:revisit-ramp-asymptotic}. Neither argument gives a work lower
bound against the complete two-phase algorithm or other local algorithms.
